\documentclass[11pt]{article}
\usepackage[margin=1in]{geometry}
\usepackage{amsmath,amssymb,amsthm,mathtools}
\usepackage{booktabs,enumitem}
\usepackage{algorithm,algpseudocode,caption}
\usepackage[hidelinks]{hyperref}

\newtheorem{theorem}{Theorem}
\newtheorem{lemma}[theorem]{Lemma}
\newtheorem{corollary}[theorem]{Corollary}
\newtheorem{proposition}[theorem]{Proposition}
\newtheorem{definition}[theorem]{Definition}
\newtheorem{assumption}[theorem]{Assumption}
\theoremstyle{remark}
\newtheorem{remark}[theorem]{Remark}

\DeclareMathOperator{\tid}{tid}\DeclareMathOperator{\pc}{pc}\DeclareMathOperator{\seq}{seq}
\DeclareMathOperator{\loc}{loc}\DeclareMathOperator{\blk}{blk}\DeclareMathOperator{\kind}{kind}
\DeclareMathOperator{\str}{str}\DeclareMathOperator{\scope}{scope}\DeclareMathOperator{\epoch}{epoch}
\DeclareMathOperator{\conf}{conf}\DeclareMathOperator{\ms}{ms}\DeclareMathOperator{\incl}{incl}
\DeclareMathOperator{\rel}{rel}\DeclareMathOperator{\acq}{acq}\DeclareMathOperator{\key}{key}
\DeclareMathOperator{\cpl}{cpl}\DeclareMathOperator{\Exec}{Exec}\DeclareMathOperator{\rf}{rf}
\DeclareMathOperator{\val}{val}\DeclareMathOperator{\Writers}{Writers}\DeclareMathOperator{\DR}{DR}
\DeclareMathOperator{\SC}{SC}
\newcommand{\PO}{<_{\mathrm{po}}}\newcommand{\HB}{\to_{\mathrm{hb}}}\newcommand{\HBr}{\trianglelefteq_{\mathrm{hb}}}
\newcommand{\CO}{\to_{\mathrm{co}}}\newcommand{\chain}{\rightsquigarrow}
\newcommand{\arr}[2]{\mathrm{arr}^{#2}_{#1}}\newcommand{\dep}[2]{\mathrm{dep}^{#2}_{#1}}\newcommand{\ext}[1]{\mathrm{exit}_{#1}}
\newcommand{\Eplus}{E^{+}}\newcommand{\Fr}{\mathcal{F}}\newcommand{\Pred}{\mathcal{P}}
\newcommand{\clk}{\mathit{clk}}\newcommand{\Bk}{B}\newcommand{\Ch}{\mathit{Ch}}\newcommand{\lastr}{\mathit{last}}
\newcommand{\arrv}{\mathit{arrived}}\newcommand{\gone}{\mathit{exited}}\newcommand{\Rep}{\mathit{Rep}}
\newcommand{\cls}[1]{\texttt{#1}}\newcommand{\M}{\mathsf{M}}
\algnewcommand{\algorithmicreport}{\textbf{report}}\algnewcommand{\Report}{\algorithmicreport{} }

\title{Happens-Before Data-Race Detection for GPU Kernels:\\ Model, Algorithm, Proofs}
\date{2026-09-26 --- against \texttt{cuVein} 3331d35; revised after the latent census (T9-0)}

\emergencystretch=2em
\begin{document}
\maketitle

\noindent This document replaces the v3 proof, the v4 delta and the 2026-09-24 pseudo-code
notes. It states one trace model, one happens-before relation (with the fence gate as a
parameter), one verdict definition, and one algorithm of which cuVein's two modes are the
two instances; proves the single-trace theorems and the per-profile certificate for it;
and lists, as corollaries, exactly where the code at 3331d35 departs from it and what
each departure costs. \emph{Sound} means no missed verdict; \emph{complete} means no
spurious verdict. \S\ref{sec:impl} was checked by hand and, for two items, on the real
oracle; the latent census (\texttt{eval/LATENT\_CENSUS.md}, 2026-09-26: the verdict layer
replayed over the kept traces of 558 programs with the detector unchanged) supplied the
corpus evidence cited in \S\ref{sec:verdict} and \S\ref{sec:impl}.

\paragraph{Status and scope (PLDI 2027).} The theorems of \S\ref{sec:single} are stated
relative to the relation of Definition~\ref{def:hb} under whichever gate the code runs:
the trusting gate at 3331d35, the instance gate of Definition~\ref{def:gate} once I4 is
implemented. The document is frozen except for the items under ``Order of work''; the
cp.async extension (I7) and the discharge of MM1--MM3 are stated as limitations, not
obligations, for the submission. Checked so far: by hand (all sections); on the real
oracle with hand-made dumps (\S\ref{sec:impl}, I1 and I2); on kept traces (the census:
\S\ref{sec:verdict} fact~(i), the latent tier, the corpus instance of $R_{\mathrm{miss}}$).
Not yet checked: the referee reading of \S\ref{sec:single} and \S\ref{sec:cert} (T6).

%=============================================================================
\section{Trace model}\label{sec:model}
%=============================================================================

\begin{definition}[Records]\label{def:records}
A trace $T$ is a finite sequence of records; $\seq(e)$ is the position of $e$. Records are:
\begin{itemize}[nosep]
\item \emph{memory records} $e$ with $\tid(e)$, $\pc(e)$, a location $\loc(e)$,
      $\kind(e)\in\{\mathtt{R},\mathtt{W},\mathtt{RMW}\}$,
      $\str(e)\in\{\mathtt{weak},\mathtt{strong}\}$ (an RMW is strong), and for strong records
      $\scope(e)\in\{\mathtt{cta},\mathtt{gpu},\mathtt{sys}\}$ with $\mathtt{cta}<\mathtt{gpu}<\mathtt{sys}$;
      $\kind$, $\str$ and $\scope$ are functions of $\pc(e)$ (the sidecar). A \emph{write} is a record with
      $\kind\ne\mathtt{R}$;
\item \emph{arrival records} $A$: $\kind(A)\in\{\mathtt{barrier},\mathtt{syncwarp}\}$, a block
      $\beta(A)$, a warp, a lane mask $m(A)$, a barrier index $\mathit{bar}(A)$, a count $n(A)$
      ($0$ = the whole block);
\item \emph{exit records} $\ext{t}$, one per thread that terminates.
\end{itemize}
A thread id is $(\beta,\omega,l)$; $\blk(t)$ is its block. A location is
$(\mathtt{shared},\beta,a)$, $(\mathtt{local},t,a)$ or $(\mathtt{global},a)$; $N_\beta$ is the
block size.
\end{definition}

\begin{definition}[Instances]\label{def:instance}
Fix a key $k=(\beta,\mathit{bar})$ and process the records of $T$ in $\seq$ order keeping
$\arrv_k$ (thread ids) and $\gone_\beta$ (threads of $\beta$ with an exit record so far). A
\texttt{barrier} arrival with key $k$ adds its lanes to $\arrv_k$; an exit adds its thread to
$\gone_\beta$. The open \emph{segment} of $k$ \emph{completes} at the first record after
which $|\arrv_k|=\exp_k$, where $\exp_k=n$ if the segment's arrivals carry a count $n>0$
and $\exp_k=N_\beta-|\gone_\beta|$ otherwise; then it forms an \emph{instance} $B$ with
participants $S(B)=\arrv_k$ and completion position $\cpl(B)$ = that record's $\seq$, and
$\arrv_k$ is emptied. A trailing incomplete segment forms no instance ($\cpl=+\infty$). A
\texttt{syncwarp} arrival is an instance by itself: $S(B)$ its masked lanes, $\cpl(B)$ its
position. (An exit can complete an instance: the last non-arrived thread of the block
leaving is what the hardware waits for.)
\end{definition}

\begin{definition}[Ghost events, program order]\label{def:po}
For each instance $B$ and $t\in S(B)$: an \emph{arrival event} $\arr{t}{B}$ at the position
of the arrival record covering $t$ and a \emph{departure event} $\dep{t}{B}$ at
$\cpl(B)^{+}$ (departures of one instance in any fixed order). $\Eplus$ is the set of memory
records, ghost events and exits, totally ordered by $\seq$. The program order $\PO$ of
thread $t$ is $\seq$ restricted to $t$'s events. The \emph{ticks} of $t$ are its arrival
events and its RMWs, in $\PO$; $\tau^k_t$ is the $k$-th; $\epoch(e)=1+\#\{$ticks of $t$
strictly $\PO$-before $e\}$, so an RMW $e$ equals $\tau^{\epoch(e)}_t$. The
\emph{processing position} of a tick is its own $\seq$ for an RMW and $\cpl(B)$ for
$\arr{t}{B}$.
\end{definition}

\begin{definition}[Well-formed trace]\label{def:wf}
$T$ is well-formed iff (W0) $\seq$ is strictly increasing; (W1) for every key, all
arrivals of a segment agree on $n$, no segment overshoots $\exp_k$, and no thread id occurs
twice in a segment; (W2) no memory, arrival or exit record of a thread lies strictly between
its arrival and the completion of that segment; (W3) no record of a thread follows its exit.
\end{definition}
W0--W3 are decidable in one pass over the records and use no clock, so the checks may
run online in the engine or offline over the dump with the same result. The
implementation's four trace-validity checks (\texttt{TV-seq-monotonic},
\texttt{TV-barrier-overfill}, \texttt{TV-barrier-completion-order},
\texttt{TV-expected-nonzero-multiwarp}) decide W0, the overshoot part of W1, and W2; the
count-agreement and duplicate-thread parts of W1, and W3, are not checked (W3 cannot be
until exit records exist). A fifth check is required once exit records exist: every open
segment is complete at the last record of the kernel ($\arrv_k=\emptyset$ for every $k$ at
the end of $T$). It is the runtime form of A3 --- a segment still open at the end means
the hardware waited for a thread the model did not, or the collector lost an arrival ---
and it would have flagged the crs-cuda kernels of I3 directly.

\begin{lemma}[$\seq$-forwardness]\label{lem:seq}
On a well-formed trace, every generator of Definition~\ref{def:hb} is $\seq$-forward; hence
$a\HB b\Rightarrow\seq(a)<\seq(b)$ and $\HB$ is a strict partial order.
\end{lemma}
\begin{proof}
(PO) by definition, using (W2) for the edge $\dep{t}{B}\PO e$. (SYNC): an arrival is at or
before $\cpl(B)$, a departure after it. (ATOM): the successor relation is $\seq$-order.
\end{proof}

\paragraph{Trusted base.} Nothing above refers to the execution. The connection is
Proposition~\ref{prop:fidelity}, under: (A1) every memory access, barrier arrival and exit
of the execution yields exactly one record with correct fields and lane mask, in each
thread's dynamic order; (A2) for every location, the $\seq$ order of its RMWs is their
coherence order; (A3) a barrier instance of Definition~\ref{def:instance} completes exactly
when the hardware barrier releases, i.e.\ the hardware waits for the non-exited threads
of the block (PTX \texttt{bar.sync} semantics) and for $n$ threads when a count is given.
Under exit records, A3 is a statement about the hardware, no longer about the collector.

%=============================================================================
\section{Happens-before and verdicts}\label{sec:hb}
%=============================================================================

\begin{definition}[Scope inclusion, moral strength]\label{def:ms}
For a strong $e$ and thread $t$: $\incl(e,t)$ iff $\scope(e)\ge\mathtt{gpu}$, or
$\scope(e)=\mathtt{cta}$ and $\blk(\tid(e))=\blk(t)$. Records $a,b$ of different threads
are \emph{morally strong}, $\ms(a,b)$, iff both are strong, $\incl(a,\tid(b))$ and
$\incl(b,\tid(a))$. (All records on a location have one size and use the generic proxy;
the monitor rejects otherwise.)
\end{definition}

\begin{definition}[Gate, chain]\label{def:gate}
A \emph{gate} is a pair of predicates $\rel,\acq$ on RMW records. It governs the (ATOM)
generator only, never $\ms$: the two-RMW exception of Definition~\ref{def:verdict} does
not depend on it. The \emph{trusting} gate sets both to $1$ everywhere. The
\emph{instance} gate uses one static predicate on the kernel CFG,
$\mathrm{fenced}(p,p',s)$: every CFG path from pc $p$ to pc $p'$ crosses a fence of
scope $\ge s$ (a \texttt{MEMBAR} of that scope; a \texttt{bar.sync} counts as a fence of
scope \texttt{cta}). For an RMW record $r$ of thread $t$ with $\PO$-previous record $p$
and $\PO$-next record $q$ (records of any kind): $\rel(r)=1$ iff $p$ does not exist or
$\mathrm{fenced}(\pc(p),\pc(r),\scope(r))$; $\acq(r)=1$ iff $q$ does not exist or
$\mathrm{fenced}(\pc(r),\pc(q),\scope(r))$. Because the paths are taken between the two
pcs the thread actually executed, a fence that lies only on the success edge of a CAS
loop qualifies for the successful CAS, a fence hoisted before a loop qualifies for the
first iteration only, and a failed CAS acquires nothing. The per-pc variant ($\acq(r)=1$
iff no memory pc is reachable from $\pc(r)$ without crossing a fence, and dually for
$\rel$) is \emph{not} used: it rejects the success-edge idiom because the failure edge is
a fenceless path in the CFG, and the latent census found exactly that on
\texttt{matrix-multiplication}'s block lock and on the ticket \texttt{atomicInc} of
\texttt{reduction}. For RMWs on one location, $q$ is the \emph{successor} of $r$ if it is
the next RMW on that location in $\seq$; $r\chain q$ iff there are RMWs
$r=z_0,\dots,z_k=q$, $k\ge1$, each the successor of the previous with
$\ms(z_i,z_{i+1})$.
\end{definition}
The instance gate admits only PTX release and acquire patterns (a fence followed in
program order by a strong write; a strong read followed by a fence) and is still
conservative --- a fence separated from the RMW by an unrelated access of the same thread
is not counted --- which is what \S\ref{sec:cert} needs. The trusting gate admits every
hand-off and is what \S\ref{sec:cert} declines to place inside PTX causality.

\begin{definition}[Happens-before]\label{def:hb}
Given a gate, $\HB$ is the least transitive relation on $\Eplus$ containing
(PO) $a\PO b\Rightarrow a\HB b$;
(SYNC) $\arr{s}{B}\HB\dep{t}{B}$ for every instance $B$ and $s,t\in S(B)$;
(ATOM) $r\HB q$ for RMWs with $r\chain q$, $\rel(r)=\acq(q)=1$.
$a\HBr b$ means $a\HB b$ or $a=b$. $\HB^{\mathrm{sync}}$ is the relation without (ATOM).
\end{definition}

\begin{definition}[Verdicts]\label{def:verdict}
For memory records with $\seq(a)<\seq(b)$: $\conf(a,b)$ iff $\loc(a)=\loc(b)$,
$\tid(a)\ne\tid(b)$ and one is a write; $\mathrm{unord}(a,b)$ iff neither $a\HB b$ nor
$b\HB a$. $\DR(a,b)$ (data race) iff $\conf\wedge\mathrm{unord}\wedge\neg\ms(a,b)$;
$\SC(a,b)$ (unordered strong conflict) iff $\conf\wedge\mathrm{unord}\wedge\ms(a,b)$ and
not both RMW. A pair is \emph{reportable} iff $\DR\vee\SC$; $\Rep(T)$ is the set of
reportable pairs with their class. Two morally strong RMWs are never reportable:
coherence orders them and atomicity makes both outcomes well-defined.
\end{definition}

\begin{lemma}[Monotonicity]\label{lem:mono}
If $\mathrel{\to_{\mathrm{hb}_1}}\subseteq\mathrel{\to_{\mathrm{hb}_2}}$ on one trace, every pair reportable under $\to_{\mathrm{hb}_2}$ is
reportable under $\to_{\mathrm{hb}_1}$ with the same class. In particular the trusting gate yields a
subset of the instance gate's reports: a sidecar that misses a fence, a scope or a strength
can add reports, never hide one. The converse does not hold for $\ms$, which the gate does
not touch: a sidecar that \emph{widens} a scope can hide a report through the two-RMW
exception, so scope is the one static label that must be conservative.
\end{lemma}
\begin{proof}
$\conf$, $\ms$, $\kind$ do not depend on $\HB$; unord is antitone.
\end{proof}

%=============================================================================
\section{The algorithm}\label{sec:alg}
%=============================================================================

One procedure, $\M\in\{\mathsf{vec},\mathsf{sync}\}$ selecting whether (ATOM) is
tracked. \emph{State}: $\clk[t]$, a clock (map thread $\to\mathbb{N}$, missing entry $0$,
$\sqcup$ pointwise max), initially $\{t\!:\!1\}$; per location $\ell$: buckets
$\Bk_\ell[u,\kappa]$ for every thread $u$ and key
$\kappa=(\kind,\str,\scope)$, each holding (epoch, pc) of $u$'s latest record on $\ell$ with
$\key(e)=\kappa$; a chain clock $\Ch_\ell$ (a clock or $\bot$) and $\lastr_\ell$, the label
$(\tid,\scope,\blk)$ of the last RMW on $\ell$; per key $k$: $\arrv_k$; per block:
$\gone_\beta$. The output $\Rep$ is a set of (record, record, class).

\begingroup\small
\noindent\rule{\textwidth}{0.6pt}
\captionof{algorithm}{$\textsc{Detect}(T,\M)$. cuVein's vector-clock mode runs
$\textsc{Detect}(T,\mathsf{vec})$ and $\textsc{Detect}(T,\mathsf{sync})$ side by side
(two clocks); scalar-clock mode runs $\textsc{Detect}(T,\mathsf{sync})$ only, whose clock
needs one shared base per instance plus one scalar per thread.}\label{alg:detect}
\noindent\rule{\textwidth}{0.4pt}
\begin{algorithmic}[1]
\Procedure{Sync}{$A$}\label{ln:sync}
  \State $J\gets\bigsqcup_{u\in A}\clk[u]$;\ \textbf{for} $u\in A$: $\clk[u]\gets J$,\ $\clk[u][u]\gets J[u]+1$
\EndProcedure
\Procedure{Complete}{$k$}\Comment{called after every arrival on $k$ and every exit in $k$'s block}
  \If{$|\arrv_k|=\exp_k$ \textbf{and} $\arrv_k\ne\emptyset$}\ \Call{Sync}{$\arrv_k$};\ $\arrv_k\gets\emptyset$\EndIf
\EndProcedure
\Procedure{Check}{$a$}\Comment{$a$ by $t$ on $\ell$; every other thread's buckets on $\ell$}\label{ln:check}
  \ForAll{$u\ne t$, $\kappa$ with $e=\Bk_\ell[u,\kappa]$ defined, $\conf(e,a)$, and $e.\epoch>\clk[t][u]$}
    \State \textbf{if} $\neg\ms(e,a)$: $\Rep\gets\Rep\cup\{(e,a,\DR)\}$;\ \textbf{else if} not both RMW: $\Rep\gets\Rep\cup\{(e,a,\SC)\}$
  \EndFor
\EndProcedure
\ForAll{records $e$ of $T$ in $\seq$ order}
  \If{$e=\ext{t}$}\ $\gone_{\blk(t)}\gets\gone_{\blk(t)}\cup\{t\}$;\ \Call{Complete}{$k$} for every open key $k$ of $\blk(t)$\label{ln:exit}
  \ElsIf{$e$ is a \texttt{barrier} arrival with key $k$}\ $\arrv_k\gets\arrv_k\cup\mathrm{lanes}(e)$;\ \Call{Complete}{$k$}\label{ln:barrier}
  \ElsIf{$e$ is a \texttt{syncwarp}}\ \Call{Sync}{masked lanes of $e$}
  \Else\Comment{memory record $a=e$ by $t$ on $\ell$}
    \If{$\M=\mathsf{vec}$ \textbf{and} $\kind(a)=\mathtt{RMW}$}\label{ln:acquire}
      \State \textbf{if} $\lastr_\ell$ defined \textbf{and} $\neg\ms(\lastr_\ell,a)$: $\Ch_\ell\gets\bot$\Comment{chain broken}
      \State \textbf{if} $\acq(a)=1$ \textbf{and} $\Ch_\ell\ne\bot$: $\clk[t]\gets\clk[t]\sqcup\Ch_\ell$\Comment{acquire}
    \EndIf
    \State \Call{Check}{$a$};\quad $\Bk_\ell[t,\key(a)]\gets(\clk[t][t],\pc(a))$\label{ln:record}
    \If{$\M=\mathsf{vec}$ \textbf{and} $\kind(a)=\mathtt{RMW}$}
      \State \textbf{if} $\rel(a)=1$: $\Ch_\ell\gets(\Ch_\ell\text{ or }\bot)\sqcup\clk[t]$\Comment{publish, with the pre-tick clock}\label{ln:publish}
      \State $\lastr_\ell\gets(t,\scope(a),\blk(t))$;\quad $\clk[t][t]\gets\clk[t][t]+1$\Comment{then tick}\label{ln:tick}
    \EndIf
  \EndIf
\EndFor
\State \Return $\Rep$
\end{algorithmic}
\noindent\rule{\textwidth}{0.6pt}
\endgroup

\medskip\noindent
Under $\M=\mathsf{sync}$ an RMW is an ordinary strong write: no chain, no acquire, no
publish, no tick, and the ticks of Definition~\ref{def:po} are arrivals only. The
trusting gate makes lines~\ref{ln:acquire}--\ref{ln:tick} unconditional in $\rel,\acq$.
The correspondence with \texttt{hb\_oracle.py}: $\clk$ = \texttt{vc} ($\mathsf{vec}$) and
\texttt{vs} ($\mathsf{sync}$), $\Ch_\ell$ = \texttt{released[loc]}, \textsc{Sync} =
\texttt{sync\_group}, \textsc{Check} = \texttt{conflict}, $\Rep$ = \texttt{hb\_races} or
\texttt{hb\_races\_sync\_only}; and \texttt{barrier\_only\_pairs} is
$\textsc{Detect}(T,\mathsf{sync})$ computed offline with $\clk[u]$ stored as a pointer to
$J$ plus the scalar $\clk[u][u]$ (after \textsc{Sync} no participant's map is ever
mutated again, so the sharing is exact). Under the instance gate the online engine does
not know $\pc(q)$ when it processes $r$; it records a pending acquire and applies the
join of line~\ref{ln:acquire} when $t$'s next record arrives, before that record is
processed. Nothing reads $\clk[t]$ in between --- other threads see $t$ only through
$\Ch_\ell$ and barrier instances, and $\clk[t][t]$ is unchanged by the join --- so the
deferred and the immediate join produce the same $\Rep$; the offline oracle, which has
the whole trace, evaluates $\acq$ directly.

%=============================================================================
\section{Single-trace correctness}\label{sec:single}
%=============================================================================

Throughout, $T$ is well-formed and $\M=\mathsf{vec}$; the $\mathsf{sync}$ case is
Corollary~\ref{cor:sync}. After the records with $\seq\le i$ are processed let
$\mathrm{last}_t(i)$ be $t$'s $\seq$-last event at a position $\le i$ (a departure if the
record at $i$ completed an instance containing $t$; an arrival if $t$ is pending), and the
\emph{frontier} $\Fr_i(t)=\{f:f\HBr\mathrm{last}_t(i)\}$. For a record $b$,
$\Pred(b)=\{f:f\HB b\}$. The immediate $\HB$-predecessors of an event are its
$\PO$-predecessor; for $\dep{t}{B}$ the arrivals $\arr{s}{B}$, $s\in S(B)$; for an RMW $q$
the RMWs $r$ with $r\chain q$, $\rel(r)=\acq(q)=1$. Hence, with $p$ the $\PO$-predecessor of
a record $b$ of $t$ and (W2),
\begin{equation}\label{eq:pred}
\Pred(b)=\Fr_{\seq(b)-1}(t)\ \cup\ \textstyle\bigcup\{\,\{f:f\HBr r\}: r\chain b,\ \rel(r)=\acq(b)=1\,\}.
\end{equation}

\begin{lemma}[Assembly]\label{lem:assembly}
\textsc{Sync} is invoked exactly once per instance $B$, at position $\cpl(B)$, with
argument $S(B)$, and never otherwise for a \texttt{barrier} key.
\end{lemma}
\begin{proof}
$\arrv_k$ and $\gone_\beta$ evolve exactly as in Definition~\ref{def:instance};
\textsc{Complete} is evaluated after every record that changes either, and (W1) makes the
count hit $\exp_k$ exactly at the segment's completing record.
\end{proof}

\begin{lemma}[Chain clock]\label{lem:chain}
After line~\ref{ln:tick} of an RMW $z$ on $\ell$,
$\Ch_\ell=\bigsqcup\{P_r: r\chain z,\ \rel(r)=1\}\sqcup(\rel(z)=1\,?\,P_z:\bot)$, where
$P_r$ is the clock published at $r$ (line~\ref{ln:publish}, before $r$'s tick).
\end{lemma}
\begin{proof}
Induction over the RMWs on $\ell$ in $\seq$ order. For the first, the chain set is empty
and the claim is line~\ref{ln:publish}. For $z$ with predecessor $p$: by IH
$\Ch_\ell=\bigsqcup\{P_r:\rel(r),\ r\chain p\text{ or }r=p\}$ before $z$'s
line~\ref{ln:acquire}. Now $r\chain z$ iff ($r\chain p$ or $r=p$) and $\ms(p,z)$; the
reset fires exactly when $\neg\ms(p,z)$; and line~\ref{ln:publish} adds $P_z$ iff
$\rel(z)$.
\end{proof}

\begin{lemma}[Clock invariant]\label{lem:vc}
After processing the records with $\seq\le i$: (a) $\clk[t][t]=1+\#\{$ticks of $t$ with
processing position $\le i\}$; (b) for $s\ne t$, $\clk[t][s]=\max(\{k:\tau^k_s\in\Fr_i(t)\}\cup\{0\})$;
(c) for every $r\chain$-maximal published RMW, $P_r[s]=\max(\{k:\tau^k_s\HBr r\}\cup\{0\})$
for $s\ne\tid(r)$ and $P_r[\tid(r)]=\epoch(r)$.
\end{lemma}
\begin{proof}
Induction on $i$; initially trivial. Two consequences of (a),(b) are used: for $s\ne t$,
$\clk[t][s]\le\#\{$processed ticks of $s\}<\clk[s][s]$, so $(\bigsqcup_{u\in S}\clk[u])[s]=\clk[s][s]$
whenever $s\in S$ (\dag), and no join raises $\clk[t][t]$. Let $e$ be the record at $i$.

\emph{Exit.} No clock changes; if it completes an instance, the next case applies.

\emph{Read or write.} No clock changes, no tick; $\Fr_i(t)=\Fr_{i-1}(t)\cup\{e\}$ by
\eqref{eq:pred} adds no tick of an $s\ne t$.

\emph{Non-completing arrival.} No clock changes; $\Fr_i(t)=\Fr_{i-1}(t)\cup\{\arr{t}{B}\}$
adds a tick of $t$ whose processing position is $\cpl(B)>i$; other frontiers unchanged.

\emph{Completing arrival or exit, instance $B$, $S=S(B)$.} For $t\in S$,
$\mathrm{last}_t(i)=\dep{t}{B}$ and $\Fr_i(t)=\{\dep{t}{B}\}\cup\bigcup_{u\in S}(\{\arr{u}{B}\}\cup\Fr_{i-1}(u))$
(by (W2), $u$ has no event between its arrival and $i$); \textsc{Sync} runs with exactly
$S$ (Lemma~\ref{lem:assembly}). For $s\notin S$ the maximal tick of $s$ in $\Fr_i(t)$ is
$\max_u\clk[u][s]=J[s]$ by (b). For $s\in S$, $\arr{s}{B}=\tau^k_s$ with $k=\clk[s][s]$
before $i$ by (a), and all earlier ticks of $s$ precede it; by (\dag) $J[s]=k$. So
$\clk[t]\gets J$ gives (b), and $\clk[t][t]\gets J[t]+1$ gives (a): exactly one tick of $t$
has processing position $i$. (c) is untouched.

\emph{RMW $q$ by $t$ on $\ell$.} By (a), $c:=\clk[t][t]=\epoch(q)$ and $q=\tau^c_t$. By
\eqref{eq:pred}, $\Fr_i(t)=\Fr_{i-1}(t)\cup\{q\}\cup\bigcup_r\{f:f\HBr r\}$ over the $r$ with
$r\chain q$, $\rel(r)=\acq(q)=1$. By Lemma~\ref{lem:chain} the acquire joins exactly
$\bigsqcup_r P_r$ over those $r$ when $\acq(q)=1$ and the chain is intact (if the chain is
broken at $p$ then $\neg\ms(p,q)$ and no $r\chain q$ exists), and by (c) $P_r[s]$ is the
maximal tick of $s$ in $\{f:f\HBr r\}$ --- for $s=\tid(r)$ this is $\epoch(r)$, as $r$ is
itself a tick of $s$ and every tick of $s$ before $r$ is $\PO$-before it. So the joined
clock satisfies (b). The clock published at $q$ has, for $s\ne t$, the maximal ticks of $s$
in $\Pred(q)=\{f:f\HBr q\}\setminus\{q\}$, and $t$-component $c=\epoch(q)$, the index of $q$
itself: this is (c) for $q$. Then $\clk[t][t]\gets c+1$ restores (a).
\end{proof}

\begin{lemma}[HB test]\label{lem:test}
Let $e$ be a record of $s$ and $a$ one of $t\ne s$ with $\seq(e)<\seq(a)$. When
\textsc{Check}($a$) runs, $\epoch(e)\le\clk[t][s]\iff e\HB a$.
\end{lemma}
\begin{proof}
At that moment $\clk[t][s]=\max\{k:\tau^k_s\in\Pred(a)\}$ by Lemma~\ref{lem:vc}(b),(c) and
\eqref{eq:pred}. ($\Leftarrow$) On a generating path from $e$ to $a$ the first edge leaving
thread $s$ is (SYNC) from an arrival or (ATOM) from an RMW of $s$, i.e.\ from a tick
$\tau^k_s$ with $e\HBr\tau^k_s$, so $\epoch(e)\le k\le\clk[t][s]$. ($\Rightarrow$) Some
$\tau^k_s\in\Pred(a)$ has $k\ge\epoch(e)$; every record of $s$ with epoch $\le k$ is
$\PO$-before-or-equal $\tau^k_s$, so $e\HBr\tau^k_s\HB a$.
\end{proof}

\begin{theorem}[Complete]\label{thm:complete}
Every triple $(e,a,c)$ added to $\Rep$ is a reportable pair of class $c$.
\end{theorem}
\begin{proof}
$e$ is in a bucket of $\loc(a)$ of a thread $u\ne t$ and $\conf(e,a)$ was tested;
$e.\epoch>\clk[t][u]$ gives $\neg(e\HB a)$ by Lemma~\ref{lem:test} and
$\seq(e)<\seq(a)$ gives $\neg(a\HB e)$ by Lemma~\ref{lem:seq}; the class is
Definition~\ref{def:verdict}'s.
\end{proof}

\begin{theorem}[Sound]\label{thm:sound}
For every reportable $(e,a)$ with $\seq(e)<\seq(a)$, \textsc{Check}($a$) adds $(e',a)$ of
the same class, where $e'=e$ or $e\PO e'$ and $\key(e')=\key(e)$.
\end{theorem}
\begin{proof}
Let $e'=\Bk_{\loc(a)}[\tid(e),\key(e)]$ when $a$ is processed: the latest record of
$u=\tid(e)$ on $\loc(a)$ with that key, so $e=e'$ or $e\PO e'$. If $e'\HB a$ then $e\HB a$,
contradiction; so $\neg(e'\HB a)$, and $\neg(a\HB e')$ by Lemma~\ref{lem:seq}. $\kind$,
$\str$, $\scope$ and the block agree between $e$ and $e'$, so $\conf$, $\ms$ and the class
agree. \textsc{Check} inspects that bucket and, by Lemma~\ref{lem:test}, reports.
\end{proof}

\begin{corollary}[The $\mathsf{sync}$ instance]\label{cor:sync}
$\textsc{Detect}(T,\mathsf{sync})$ satisfies Lemmas~\ref{lem:vc} and~\ref{lem:test} and
Theorems~\ref{thm:complete} and~\ref{thm:sound} with $\HB^{\mathrm{sync}}$ in place of
$\HB$ and ticks = arrivals. Hence a pair is $\HB^{\mathrm{sync}}$-ordered iff it is
ordered by barrier and syncwarp joins alone, which are schedule-independent.
\end{corollary}
\begin{proof}
Delete the RMW case of Lemma~\ref{lem:vc} and the (ATOM) term of \eqref{eq:pred}; nothing
else in the proofs mentions them.
\end{proof}

\begin{remark}[Why one bucket per key]\label{rem:buckets}
With a single ``last write'' per location the theorem fails: $w_0$ strong by $A$, $w_1$
strong by $B$, unordered (so no report: $\ms$, both writes, not both RMW would give $\SC$ ---
which the implementation drops, \S\ref{sec:impl}); a barrier joins $B$ and $C$; a weak read
$r$ by $C$. $(w_0,r)$ is a $\DR$; a check against $w_1$ alone sees $w_1\HB r$ and reports
nothing. A bucket representative is always $\PO$-related to the record it replaced, which
is what the proof of Theorem~\ref{thm:sound} uses.
\end{remark}

%=============================================================================
\section{The verdict layer}\label{sec:verdict}
%=============================================================================

$\Rep$ is a set of pairs of records; cuVein reports per \emph{pc pair}
$\{u,v\}$ and asks a second question: is the pair ordered in \emph{every} schedule? A
static certificate $\textsc{Static}(u,v,d)$, from the CFG region graph and the observed
thread distance $d$, returns $\rho$ (R1: every path between the two regions crosses a
qualifying barrier of scope $\ge d$, with loops handled by cutting; R2: both pcs coherent
at scope $\ge d$ --- the $\ms$ exception in certificate form, not an ordering) and $\chi$
(R3: a chain of atomic pcs from $u$ to $v$, each hop an observed atomic--atomic edge at
scope $\ge d$, release-fenced at the start, ending in a pc that dominates $v$ and whose
critical section $v$ has not left).

Two amendments to R3 follow from the census. Its \emph{release point} is taken from the
trace, not from region dominance: the first RMW $r$ of the chain is the $\PO$-next RMW of
$u$'s thread after $u$, and ``release-fenced'' is $\mathrm{fenced}(\pc(u),\pc(r),d)$ of
Definition~\ref{def:gate}. The region test declined 23 of the census's 25 race-free
latent pairs (\texttt{rule-110}, \texttt{reduction}) because the protected access sat in
a loop or branch whose region dominates no atomic, although the release followed it in
program order and 21 of the hand-offs were fully fenced; the trace has the program order
exactly, and only the fence question is static. And R3 is made symmetric in the two
accesses: it declines when $v$ has left its critical section (\texttt{\_past\_release},
which is what catches \texttt{rtraw}) and must equally decline when $u$ has not entered
one, i.e.\ when $u$ precedes in its own thread an acquire on the chain's location
(write-before-lock, the mirror of rtraw; not exercised by the census, to be checked in
\texttt{HBGraph.chain}). With $\Rep_{\mathsf{vec}}$ and $\Rep_{\mathsf{sync}}$ projected
to pc pairs, the verdict for a candidate pair is
\[
\begin{array}{ll}
\{u,v\}\in\Rep_{\mathsf{vec}}\text{ with class }\SC: & \textsc{Strong conflict},\ \cls{sc}\\
\{u,v\}\in\Rep_{\mathsf{vec}}\text{ with class }\DR: & \textsc{Race},\ \text{class \cls{model\_bug} if }\rho\text{ else \cls{structural}}\\
\rho\vee\chi: & \textsc{Ordered},\ \cls{ordered}\\
\{u,v\}\notin\Rep_{\mathsf{sync}}: & \textsc{Ordered},\ \cls{barrier-ordered}\\
\text{otherwise}: & \textsc{Latent},\ \cls{latent}
\end{array}
\]
and scalar-clock mode is the same without the first two rows and with \textsc{Race} in
the last: it has no $\Rep_{\mathsf{vec}}$ and cannot tell a race of this run from a latent
pair. The first row exists because R2 certifies exactly the morally strong pairs: once
I2 reports $\SC$, every $\SC$ pair would otherwise land in \cls{model\_bug}. \textsc{Latent}
is a third verdict, not a race: the pair is ordered in this execution, necessarily
through an atomic hand-off, and no schedule-independent evidence orders it. It is
reported, counted as neither true nor false positive in the precision tables, and
evaluated on its own (D8, below).

Four facts follow from \S\ref{sec:single}. (i) The full clock never grants
\textsc{Ordered}: it vetoes a certificate (row two) or separates \textsc{Race} from
\textsc{Latent}, so scalar-clock's \textsc{Race} set is vector-clock's
$\textsc{Race}\cup\textsc{Latent}$ minus the vetoes. The census confirmed this on the same
trace for 557 of 558 programs; the exception is \texttt{crs-cuda}, where the scalar-clock
offline pass is skipped above \texttt{CUVEIN\_BARRIER\_PASS\_MAX\_LANES} (26 kernels above
$5\cdot10^{6}$ lane-accesses) and six vector-clock pairs are never judged. (ii)
\cls{barrier-ordered} is justified by Corollary~\ref{cor:sync}: absence from
$\Rep_{\mathsf{sync}}$ is schedule-independent ordering. (iii) \textsc{Race} in
vector-clock mode is exactly $\Rep_{\mathsf{vec}}$ projected to pc pairs, so
Theorems~\ref{thm:complete} and~\ref{thm:sound} apply to that verdict as stated and
Theorem~\ref{thm:cert} lifts it to every execution with the same profile; no theorem
speaks about \textsc{Latent}. (iv) \cls{model\_bug} is the disagreement of two
evaluations of the same (SYNC) generator: on a faithful trace, if R1 certifies $\{u,v\}$
then every dynamic instance of the pair is $\HB^{\mathrm{sync}}$-ordered (each thread of
the block passes the certified barrier between the two regions, and block-wide instances
are totally ordered), so the cell is empty; it has fired only on defects (F2, F6, and
crs-cuda's exit-unaware assembly, all fixed or covered by \S\ref{sec:impl}).

\paragraph{What scalar-clock mode misses.} By Lemma~\ref{lem:mono},
$\Rep_{\mathsf{sync}}\supseteq\Rep_{\mathsf{vec}}$ on every trace: scalar-clock mode
misses no record-level race. Its misses are in the verdict layer, and they are exactly
three. (a) R3 grants \textsc{Ordered} at pc level to a pair some dynamic instance of
which is unordered --- the $\exists$-as-$\forall$ flaw of pc granularity; vector-clock
mode vetoes it through row two, scalar-clock cannot. The P5 canary is the corpus
instance, and the only program of the census's 558 whose verdict differs between the
modes. (b) Candidate generation: a conflicting pair with no dependency edge and no
event-stream candidate is never judged, whereas the engine checks every location. (c) The
\texttt{CUVEIN\_BARRIER\_PASS\_MAX\_LANES} cutoff, under which the offline pass and its
event-stream candidates are skipped (crs-cuda). Everything else the two modes report
identically.

\paragraph{The latent tier (D8, resolved).} \textsc{Ordered} keeps its meaning ``in
every schedule''; \textsc{Race} means ``unordered in this execution''; what lies between
is \textsc{Latent}. The census measured the tier on the detector at 3331d35: it alone
reports 10 of the 236 labelled races vector-clock mode finds, all ScoR litmus kernels,
and 25 pairs on 5 race-free ScoR applications. Of the 10, eight are fence and scope bugs
that the run did \emph{not} order under PTX --- the trusting gate (I4) joined the hand-off
regardless, and R3's fence check is what kept the verdict; one
(\texttt{race\_interblock\_fence\_rtraw}) is a race of the recorded run missed through I1,
the first corpus instance of $R_{\mathrm{miss}}$; one
(\texttt{race\_interblock\_none-lock\_rtraw}) is genuinely ordered by the recorded
acquisition order and is a race only under the other one. Of the 25, 23 are the R3
release-point defect above and 2 the success-edge idiom that the per-pc gate rejects. So
after I1, I4 and the R3 amendments, the tier's own content is the rtraw kind: a pair
ordered in this run only through a fenced lock hand-off whose order the program does not
force. That class is real (neither racecheck nor iGUARD reports it; the other baselines
were not run on P5), and it cannot be given a witness from the trace: values are not
recorded, so a reordering of the critical sections cannot be checked to preserve
reads-from, which is what a predictive analysis would need. It is therefore reported as a
warning with its own precision, confirmed where possible by re-recording in the other
acquisition order, and the paper's numbers are given at two operating points,
\textsc{Race} alone and $\textsc{Race}\cup\textsc{Latent}$ (the latter is what today's
tables count).

%=============================================================================
\section{From one trace to all executions}\label{sec:cert}
%=============================================================================

\begin{proposition}[Fidelity]\label{prop:fidelity}
Let $T$ be a well-formed trace of execution $E$ and $\to_{\mathrm{hb},E}$ the relation of
Definition~\ref{def:hb} over $E$'s own events, barrier instances and coherence orders.
Under A1--A3, $\to_{\mathrm{hb},T}=\to_{\mathrm{hb},E}$ on memory records and $\Rep(T)=\Rep(E)$.
\end{proposition}
\begin{proof}
(PO) by A1. (SYNC): by A1 each arrival and exit is recorded; by A3 an instance completes
exactly when the hardware releases, so $S(B)$ is the true participant set. (ATOM): successor
and $\ms$ depend only on coherence order (A2) and static labels.
\end{proof}

The single-trace theorems use nothing about the platform. The certificate below is
relative to three memory-model obligations, to be discharged against the PTX
formalisation (Lustig et al., ASPLOS 2019) with the instance gate: (MM1) $\CO\cup\rf$ is
acyclic, where $\CO$ is $\HB$ plus the completion edges $\ext{s}\CO\dep{t}{B}$ for every
non-participant $s$ of $B$'s block; (MM2) coherence: a read or RMW $e$ does not read from a
$w$ with $w\HB w'\HB e$ for some write $w'$, and reads from a write if some write
$\HB$-precedes it; (MM3) every thread of a block either participates in each of the block's
instances or exits, and executions terminate. Under the instance gate every (ATOM)
edge, extended by (PO) on both sides, is a release/acquire pattern on a morally strong
chain, hence inside PTX causality ($\mathrm{fenced}$ on the thread's own path supplies the
release and acquire patterns, $\ms$ along the chain the observation order); under the
trusting gate this containment is not claimed, and the certificate is then a statement
about the implemented, stronger $\HB$. The containment is argued in one paragraph once
the gate is implemented (I4); it is not a separate proof obligation for the submission.

Fix $P$ deterministic, $I$, and a \emph{profile} $\Pi$ (for each location, the sequence
of (thread, per-thread index) of its RMWs in coherence order); $\Exec(P,I,\Pi)$ are the
executions with that profile. Label every event by its thread and index in the thread's
order, and instances by (block, count); the \emph{counterpart} $\varphi(e')$ of an event of
$T'$ is the event of $T$ with the same label.

\begin{lemma}[First reportable prefix]\label{lem:first}
If $\Rep(T')\ne\emptyset$, let $L'$ linearise $\CO'\cup\rf'$ (MM1), $b$ the $L'$-least later
member of a reportable pair, $Q=\{e:e<_{L'}b\}$. Then $Q$ contains no reportable pair and
is closed under $\CO'$-predecessors.
\end{lemma}

\begin{lemma}[Matching]\label{lem:match}
Let $E\in\Exec(P,I,\Pi)$ have a well-formed trace $T$ with $\Rep(T)=\emptyset$, and
$E'\in\Exec(P,I,\Pi)$ a well-formed $T'$. With $Q,b$ as above (or $Q=\Eplus(T')$ if
$\Rep(T')=\emptyset$), every $e\in Q\cup\{b\}$ has a counterpart agreeing in kind, pc,
location, labels and written value; if $e\in Q$ is a read or RMW, $\val_r(e)=\val_r(\varphi e)$;
$g\in Q$, $g\HB' e\Rightarrow\varphi g\HB\varphi e$; and $g\HB\varphi e\Rightarrow g=\varphi g'$
for some $g'\in Q$ with $g'\HB' e$.
\end{lemma}
\begin{proof}
Induction along $L'$; let $t=\tid(e)$. \emph{Participants.} If $\dep{u}{B}\HBr' e$ for an
instance labelled $(\beta,k)$, then $S'(B)=S(B_{(\beta,k)})$: a participant $s$ of $B$ has
all events up to its $k$-th arrival matched (IH), so it reaches that barrier in $E$; a
participant of $B_{(\beta,k)}$ absent from $S'(B)$ has, by MM3, an exit in $T'$ with
$\ext{s}\CO'\dep{u}{B}\HBr' e$, so all its events are in $Q$ and matched with the same read
values, hence the same state --- and a deterministic thread cannot both exit and continue
from one state. \emph{Existence.} $e$'s $\PO$-predecessors are in $Q$ and matched with equal
read values; determinism fixes the next instruction, address and written value, and $T$
contains it since $t$ has not exited in $T$. \emph{Forward.} Each generator maps: (PO) by
labels, (SYNC) by the participant step, (ATOM) because equal profiles make $\varphi g$ the
coherence predecessor of $\varphi e$ with the same labels and blocks. \emph{Backward.} An
immediate predecessor of $\varphi e$ is the image of the corresponding $\PO$-predecessor,
arrival (participant step) or coherence predecessor (equal profiles) of $e$, each in $Q$;
chains compose by IH. \emph{Values.} Let $w'=\rf'(e)$, $W=\rf(\varphi e)$. In $T$, no
reportable pair and MM1 make the writers of $\varphi e$ an $\HB$-chain --- two writes on a
location not both RMW are $\HB$-comparable, and RMW order is fixed by $\Pi$ --- whose
maximum is $W$ by MM2; in $T'$, $w'\in Q$, $(w',e)$ not reportable and $L'\supseteq\HB'$ give
$w'\HB' e$ and, by MM2, no writer of $e$ above $w'$. If $\varphi w'\ne W$ then
$\varphi w'\HB W\HB\varphi e$; backward preservation at $e$ and at the preimage of $W$
yields $w'\HB' W'\HB' e$ with $W'$ a writer of $e$, contradicting MM2. So $W=\varphi w'$ and
the read values agree (the init case is the same with $\Writers=\emptyset$).
\end{proof}

\begin{theorem}[Per-profile certificate]\label{thm:cert}
If some $E\in\Exec(P,I,\Pi)$ has a well-formed trace with $\Rep(T)=\emptyset$, then every
$E'\in\Exec(P,I,\Pi)$ has $\Rep(T')=\emptyset$, and $T'$ consists of the same events as $T$
up to $\seq$ (same labels, values, instances, participants, hence the same $\HB$).
\end{theorem}
\begin{proof}
Suppose $\Rep(T')\ne\emptyset$ and take $a,b,Q$ from Lemma~\ref{lem:first}. By
Lemma~\ref{lem:match}, $\varphi a,\varphi b$ exist and conflict with the same labels. If
$\varphi a\HB\varphi b$, backward preservation at $b$ gives $a'\in Q$ with $\varphi a'=\varphi a$
and $a'\HB' b$, and labels force $a'=a$, contradicting the pair; if $\varphi b\HB\varphi a$,
backward preservation at $a$ puts $b\in Q$, contradiction. So $(\varphi a,\varphi b)\in\Rep(T)$,
contradiction. With $Q=\Eplus(T')$, Lemma~\ref{lem:match} gives a label-preserving
bijection that preserves and reflects $\HB$.
\end{proof}

\begin{corollary}\label{cor:causal}
(a) If some execution in $\Exec(P,I,\Pi)$ has a reportable pair, every execution in it has
one, and by Theorem~\ref{thm:sound} the algorithm reports on any of their well-formed
traces. (b) If $P$ has no RMW, $\Pi$ is empty and one clean trace certifies $P(I)$ under all
schedules. (c) A trace with no $\DR$ but with $\SC$ pairs certifies nothing: the coherence
order of unordered strong loads and stores is a schedule input the profile does not fix.
(d) With a per-thread random outcome $R$ fixed, everything holds per $(R,\Pi)$; across
different $R$ nothing transfers.
\end{corollary}

%=============================================================================
\section{The implementation at 3331d35}\label{sec:impl}
%=============================================================================

\texttt{hb\_oracle.py} (= \texttt{HbEngine}), \texttt{barrier\_only\_pairs} and the verdict
layer of \texttt{sync\_dominance.py} implement \S\ref{sec:alg}--\ref{sec:verdict} with the
substitutions below. Each is stated with what it costs against \S\ref{sec:single}.

\begin{center}\small
\begin{tabular}{@{}lp{6.4cm}p{6.2cm}@{}}
\toprule
& implementation & consequence \\
\midrule
I1 & \emph{tick before publish}: line~\ref{ln:tick} before line~\ref{ln:publish}; the RMW's bucket entry carries the post-tick epoch & Theorem~\ref{thm:sound} fails on $R_{\mathrm{miss}}$ (Proposition~\ref{prop:rmiss}); Theorem~\ref{thm:complete} holds; Corollary~\ref{cor:causal}(a) void. Corpus instance: \texttt{race\_interblock\_fence\_rtraw} (evcand vector-clock dump; the verdict survived only as \cls{latent}) \\
I2 & \emph{one last write and one read per thread} in place of buckets; $\SC$ dropped (\texttt{--strong-ldst generic}: coherent pairs never conflict) & Theorem~\ref{thm:sound} fails for $\DR$ (Remark~\ref{rem:buckets}); with policy \texttt{none} (every strong load/store weak) the FastTrack argument applies and Theorem~\ref{thm:sound} holds under Definition~\ref{def:verdict} with $\ms$ restricted to RMW pairs \\
I3 & \emph{exits ignored}: no exit records in the dump, $\exp_k=N_\beta$; no end-of-kernel completeness check & A3 false for early-exit kernels; instances never complete; every store--barrier--load pair reported (crs-cuda: 155/7 spurious) and, absent the fifth check of \S\ref{sec:model}, unflagged \\
I4 & \emph{trusting gate} & $\HB$ is a superset of the instance-gated one; Lemma~\ref{lem:mono}: reports are a subset; the certificate is not inside PTX causality (\texttt{test\_relaxed\_handoff\_should\_race}, strict xfail). Corpus: 8 of the 10 latent-only ScoR races (\S\ref{sec:verdict}) are unordered under PTX and ordered by this gate; only R3's fence check kept their verdicts. Implement the instance gate of Definition~\ref{def:gate}, not the per-pc one \\
I5 & local memory keyed by $(\mathtt{local},a)$, not per thread & spurious WAW if the collector records the per-thread window offset; to check \\
I6 & monitor: row ``re-arriving warp'' per warp, not per lane; per-warp degrade when the count is unknown & over-rejects under independent thread scheduling (predicted from the code, not observed: 0 violations on ScoR 32/32, \texttt{HARDENING\_REPORT.md}); the degrade path is unreachable on launched kernels and unsound without the monitor, so it is deleted rather than kept \\
I7 & cp.async copies as accesses of a virtual agent $t\,|\,2^{62}$ with commit/wait joins & outside the model; needs generators (ISSUE), (WAIT) and a re-proof of \S\ref{sec:single} \\
\bottomrule
\end{tabular}
\end{center}

\begin{proposition}[Missed class of I1]\label{prop:rmiss}
With I1 alone, the reports are a subset of $\textsc{Detect}$'s and the omitted pairs are
exactly $R_{\mathrm{miss}}=\{(a,b): a$ a non-RMW record of $t$ whose $\PO$-last preceding
tick is an RMW $r$ with $r\HBr b$, no later tick of $t$ in $\Pred(b)$, $\seq(a)<\seq(b)\}$.
\end{proposition}
\begin{proof}
Let $c=\epoch(r)$. I1 publishes $\{t:c{+}1\}$ where line~\ref{ln:publish} publishes
$\{t:c\}$; arrivals publish the same in both; joins propagate. So for $u\ne t$,
$\clk_{\mathrm{I1}}[u][t]=\clk[u][t]+\delta$ with $\delta=1$ iff the maximal tick of $t$ in
$u$'s frontier is an RMW reached through (ATOM). For a non-RMW $a$ of $t$ and $\delta=1$
the tests $\epoch(a)>c$ and $\epoch(a)>c{+}1$ differ exactly at $\epoch(a)=c{+}1$: the
records of $t$ between $r$ and its next tick, unordered with $b$ --- $R_{\mathrm{miss}}$.
For an RMW $a$ at tick $c'$ the test $c'{+}1>\clk+\delta$ equals $c'>\clk$ in both cases
($\delta=0$ forces $\clk\ne c'$). In the opposite $\seq$ order $a$'s own check uses $t$'s
clock of $u$, which is unaffected.
\end{proof}
\noindent\emph{Tested} (\texttt{check\_e1\_e2.py}, real oracle, hand-made dumps): the
\texttt{rtraw} lock pattern with block~0 first --- I1 reports nothing, the reordered oracle
reports the RAW; the other schedule is reported by both. Remark~\ref{rem:buckets}'s kernel:
policy \texttt{generic} reports nothing (in both clocks), policy \texttt{none} reports
$(w_0,w_1)$.

\noindent\emph{Observed on a kept trace} (latent census). The evcand vector-clock dump of
the ScoR kernel \texttt{race\_interblock\_fence\_rtraw} has block~0 releasing at seq~5
with \texttt{atomicExch(\&flag,1)}, block~1's spin succeeding at seq~6, block~0 reading
\texttt{data} at seq~7 after its own release and block~1 writing it at seq~8, with
\texttt{hb\_races} empty --- $R_{\mathrm{miss}}$ on a real trace, proved-in-effect on
this one dump, not re-run with a publish-then-tick oracle. \emph{Monitor}: one unit test
per check that must raise plus a valid multi-warp control that must not (7/7 in
\texttt{test\_barrier\_soundness.py}); zero violations on the ScoR corpus; the
head-to-head suites are not yet counted.

\paragraph{Order of work (PLDI 2027).}
(1) D8 as resolved in \S\ref{sec:verdict}: three verdicts, two operating points, the
$\SC$ row. (2) I3: exit records in the dump, \textsc{Complete} on exit, the fifth
monitor check of \S\ref{sec:model}; A3 becomes PTX semantics (D7, route (a)). (3) I1, I2
(buckets, D6) and I5 in one change to oracle and engine, with the offline re-score;
\texttt{race\_interblock\_fence\_rtraw} is the corpus regression test for I1, litmus
O6(vi) for I2. (4) Sidecar strength and scope (O1): \texttt{.STRONG} loads and stores
and \texttt{volatile} classed strong with their scope, atomicity no longer decided by
the RMW opcode alone --- the DR/SC split is wrong in both modes until then, and the
scope label must err narrow (Lemma~\ref{lem:mono}, converse). (5) I4 as the instance
gate of Definition~\ref{def:gate}: the fence inventory per SM (O2) and the
$\mathrm{fenced}$ predicate from the CFG dots (O3), shared with R3's release-point test;
the engine defers the acquire join (\S\ref{sec:alg}). The gate is measured before it is
wired: the census's fence classification is the before; the expected after is the eight
ScoR fence races moving from \textsc{Latent} to \textsc{Race}, \texttt{matrix-multiplication}
staying clean, and the ticket idiom of \texttt{reduction} reported (a PTX race by the
letter, which iGUARD also reports; a labelled-clean program, to be footnoted, not
suppressed). (6) The vector clock's memory (T5b): the reference mode has no dump for the
realistic suites, which bounds the mode comparison to litmus and micro-benchmarks until
it runs there. Deferred past the submission: I6, I7 (model and re-proof; stated as a
limitation), the discharge of MM1--MM3. Obligations kept from v4: uniform access size and
generic proxy per location (monitor), and ``if some RMW on $\ell$ releases, every write
on $\ell$ is an RMW'' (else clear $\rel$ on $\ell$; conservative by
Lemma~\ref{lem:mono}).

\end{document}
